\section{Handling confidential information}

\begin{frame}{Problem statement}
  \begin{itemize}
    \item Most systems handle some type of confidential information
    \begin{itemize}
      \item Customer/User data
      \item Intellectual Property
      \item Cryptographic material (private keys)
    \end{itemize}
    \item The usual way to protect it is \textbf{encryption}
    \item Raises the issue of storing the \textbf{key}
    \item This is a hard problem on embedded systems
    \begin{itemize}
      \item adversaries might have physical access to the system
      \item unattended boot is incompatible with keys derived from user input
    \end{itemize}
  \end{itemize}
\end{frame}

\subsection{Hardware Security Modules (HSMs)}

\begin{frame}{HSMs}
  \begin{itemize}
    \item Dedicated Hardware for storing cryptographic secrets
    \item Implements cryptographic operations:
    \begin{itemize}
      \item key generation
      \item encryption
      \item signing
    \end{itemize}
    \item Prevent the key from leaving in cleartext form
    \item Some can be exported in \textbf{wrapped} (encrypted) form
    \item Separate usage of various keys based on PINs
  \end{itemize}
\end{frame}

\begin{frame}{HSMs: threat model}
  \begin{itemize}
    \item HSMs split 
    \begin{itemize}
      \item \textbf{usage} of cryptographic material from
      \item \textbf{knowledge} of the cryptographic material
    \end{itemize}
    \item In case of compromise, this means recovery can happen without rotation
    \item HSMs do not necessarily prevent an adversary from \textbf{using} the key
  \end{itemize}
\end{frame}

\begin{frame}{HSMs: usage}
  \begin{itemize}
    \item HSMs are usually used to store sensitive keys
    \begin{itemize}
      \item Secure boot root keys
      \item PKI root CA private keys
      \item DNSSEC signature keys
      \item Blockchain account private keys (cryptographic wallets)
    \end{itemize}
    \item Some HSMs have an on-device interface (screen + buttons)
    \item All HSMs implement a command interface, and a PKCS\#11 API
  \end{itemize}
\end{frame}

\begin{frame}{"Software HSMs"}
  \begin{itemize}
    \item Strange concept at first glance
    \item This is the difference between "hot" and "cold" crypto wallets
    \item Much less secure than an actual hardware device
    \item Also less expensive, and potentially easier to use
    \item Examples: 
    \begin{itemize}
      \item \href{https://www.softhsm.org/}{SoftHSM}
      \item OP-TEE has a \href{https://github.com/OP-TEE/optee_os/tree/master/ta/pkcs11}{PKCS\#11 TA}
    \end{itemize}
  \end{itemize}
\end{frame}

%\begin{frame}{HSMs: Examples}
%\end{frame}

\subsection{PKCS\#11}

\begin{frame}{PKCS}
  \begin{itemize}
    \item Public Key Cryptography Standards
    \item Set of standards published by RSA Security
    \item Some are more relevant than others:
    \begin{itemize}
      \item PKCS\#1, or \href{RFC 8017}{https://datatracker.ietf.org/doc/html/rfc8017}
            is the RSA specification
      \item PKCS\#3 describes the Diffie-Hellman key agreement protocol
      \item PKCS\#7 or \href{RFC 2315}{https://datatracker.ietf.org/doc/html/rfc2315}
            describes the Cryptographic Message Syntax (CMS)
      \item PKCS\#10 is a common format for Certificate Signing Requests (CSRs)
      \item PKCS\#11, or "Cryptoki" is the Cryptographic Token Interface
    \end{itemize}
  \end{itemize}
\end{frame}

\begin{frame}{PKCS\#11}
  \begin{itemize}
    \item CRYPtographic TOKen Interface
    \item v1.0 published by RSA Security in 1995
    \item Libraries implementing it are often called cryptoki
    \item Since 2013, overseen by an \href{https://www.oasis-open.org/committees/tc_home.php?wg_abbrev=pkcs11}{OASIS technical committee}
    \item The specification includes C header files
    \item It is then up to token manufacturers to implement the API
    \begin{itemize}
      \item Nitrokeys' \href{https://github.com/Nitrokey/nethsm-pkcs11}{nethsm-pkcs11}
      \item Yubico's \href{https://developers.yubico.com/yubihsm-shell/yubihsm-pkcs11.html}{PKCS11 module for YubiHSM}
      \item Thales' libCryptoki2, part of the Luna client
    \end{itemize}
  \end{itemize}
\end{frame}

\begin{frame}{PKCS\#11 clients}
  \begin{itemize}
    \item PKCS\#11 is very specific (even implements a header)
    \item Token-specific implementations are usually shared objects
    \item So we can have a generic client compatible with all modules:
    \begin{itemize}
      \item OpenSSL, via a \href{https://github.com/openssl-projects/pkcs11-provider}{provider}
      \item GnuTLS' \href{https://www.gnutls.org/manual/html_node/p11tool-Invocation.html}{p11tool}
    \end{itemize}
  \end{itemize}
\end{frame}

\subsection{Cryptographic keys in the kernel}

\begin{frame}{The Kernel Key Retention Service}
  \begin{itemize}
    \item Also called "Kernel Key Ring Service", or "Kernel keyring"
    \item Documented in \kdochtml{security/keys/core}
    \item Can essentially act as a software HSM
    \item The protection here is against adversaries that have gotten
          unprivileged code execution
  \end{itemize}
\end{frame}

\begin{frame}{Trusted \& Encrypted Keys}
  \begin{itemize}
    \item Specific key types in the Kernel keyring
    \item Stored in plaintext in kernel memory, only exportable wrapped
    \item The kernel sort of acts as an HSM for userland
    \item \textbf{Trusted key} encryption is harware-backed
    \item \textbf{Encrypted key} encryption is backed by a kernel keyring key
  \end{itemize}
\end{frame}
